\documentclass{article}
\usepackage{amsmath} % Para ecuaciones
\usepackage{amssymb}
\usepackage{graphicx} % Para gráficos
\usepackage{url} % Para URLs de referencias

\title{Sistema de Hormigas Elitista}
\author{Jesús Samuel García Carballo}
\date{\today}

\begin{document}

\maketitle

\section{Introducción}

Uno de los problemas más estudiados por los etólogos ha sido entender cómo animales con capacidades visuales
limitadas, como las hormigas, logran establecer rutas eficientes desde su colonia para conseguir recursos y volver.
Se ha descubierto que la forma más usual para comunicar información entre los individuos sobre los caminos, 
y que se usa para decidir a dónde ir, se basa en el rastro de feromonas.
El movimiento de hormigas deposita esta sustancia en el suelo, marcando el sendero.
Este mecanismo genera un proceso de retroalimentación positiva: de tal forma que cuanto más sigan las hormigas 
un rastro, más atractivo se vuelve ese rastro para ser seguido. 
En concreto, este comportamiento colectivo se caracteriza porque la probabilidad de que una hormiga escoja 
un camino aumenta en función del número de individuos que ya han transitado esa misma ruta.

\section{Sistema de Hormigas}

En primer lugar vamos a utilizar el problema del viajante de comercio (TSP) como caso de estudio para
implementar nuestro sistema de hormigas elitista. 
El TSP es un problema clásico de optimización combinatoria que busca encontrar la ruta más corta 
que visite un conjunto de ciudades exactamente una vez y regrese a la ciudad de origen.

Dado un conjunto de n ciudades, denominamos $d_{ij}$ a la longitud del camino entre la ciudades i y j,
que suele ser la distancia euclídea. Este conjunto de ciudades y caminos viene representado por un grafo (N, E),
donde N es el conjunto de nodos (ciudades) y E el conjunto de aristas (caminos entre ciudades).

Denominamos $b_i(t)$ el número total de hormigas en la ciudad i en el tiempo t y $m = \sum_{i=1}^{n} b_i(t)$ 
el número total de hormigas en el tiempo t.
Cada hormiga dispone de las siguientes características:
\begin{itemize}
    \item Escoge la ciudad siguiente con una probabilidad, que es función de la cantidad de feromona en el camino y de la distancia a la ciudad siguiente.
    \item Obligamos a que realice recorridos válidos, prohibiendo las transiciones a pueblos ya visitados hasta que se complete un recorrido, mediante una lista tabú.
    \item Cuando se completa un recorrido, deposita una sustancia llamada rastro en cada arista visitada.
\end{itemize}

Sea $\tau_{ij}(t)$ la cantidad de feromona en el camino entre las ciudades i y j en el tiempo t.
Cada hormiga en el tiempo t, debe escoger la siguiente ciudad, donde estará en tiempo t+1.
Por lo tanto, si llamamos a una iteración del algoritmo a los m movimientos realizados por las m hormigas en el intervalo (t, t+1), 
entonces cada n iteraciones del algoritmo, cada hormiga ha completado un recorrido, lo que denominaremos ciclo.

En cada ciclo la intensidad del camino se actualiza de la siguiente forma:
\begin{equation}
    \tau_{ij}(t+n) = \rho \cdot \tau_{ij}(t) + \Delta \tau_{ij}
\end{equation}
donde $0 < \rho < 1$ es un coeficiente tal que $1 - \rho$ representa la evaporación del camino entre t y t+n, y
\begin{equation}
   \Delta \tau_{ij} = \sum_{k=1}^{m} \Delta \tau_{ij}^k
\end{equation}
donde $\Delta \tau_{ij}^k$ es la cantidad por unidad de longitud de feromona que la hormiga k deposita en el camino entre las ciudades i y j, 
entre el tiempo t y t+n.
Esta cantidad se define como:
\begin{equation}
    \Delta \tau_{ij}^k = \begin{cases}
    Q / L_k & \text{si la hormiga } k \text{ utiliza el camino } (i, j) \text{ en su recorrido}\\
    0 & \text{en otro caso}
    \end{cases}
\end{equation}
donde Q es una constante y $L_k$ la longitud del recorrido realizado por la hormiga k.

Con el objetivo de satisfacer la restricción de que cada hormiga debe visitar todas las n diferentes ciudades, asociamos a cada hormiga
una lista tabú que almacena las ciudades ya visitadas por esa hormiga.
Cuando se completa un ciclo, la lista tabú se utiliza para calcular la solución actual de cada hormiga.
Luego, la lista se vacía y cada hormiga comienza un nuevo ciclo.

Llamamos visibilidad, $n_{ij}$, a la distancia inversa entre las ciudades i y j, $\frac{1}{d_{i,j}}$. 
Esta cantidad no se modifica durante la ejecución del algoritmo y representa la información heurística disponible para las hormigas.

Definimos la probabilidad de transición de una ciudad i a una ciudad j para la hormiga k en el tiempo t como:
\begin{equation}
    p_{ij}^k = \begin{cases}
    \frac{[\tau_{ij}(t)]^{\alpha} \cdot [n_{ij}]^{\beta}}{\sum_{l \in \text{allowed}_k} [\tau_{il}(t)]^{\alpha} \cdot [n_{il}]^{\beta}} & \text{si } j \in \text{allowed}_k\\
    0 & \text{en otro caso}
    \end{cases}
    \label{eq}
\end{equation}
donde $\text{allowed}_k$ es el conjunto de ciudades que la hormiga k no ha visitado aún, y $\alpha$ y $\beta$ son parámetros que controlan la
importancia relativa entre la feromona y la visibilidad.
Por lo tanto, esta probabilidad es un equilibrio entre la visibilidad, que indica que se deben elegir las ciudades más cercanas con mayor probabilidad,
y la intensidad del rastro en el tiempo t, que indica que si en la arista (i, j) ha habido mucho tráfico, entonces ese camino es más atractivo.

\section{Algoritmo}

El algoritmo comienza con una fase de inicialización, donde las hormigas se posicionan en diferentes ciudades y se establece la intensidad inicial de feromona en cada arista del grafo
mediante $\tau_{ij}(0)$. Inicializamos el primer elemento de la lista tabú de cada hormiga con la ciudad donde se encuentra.

Luego, cada hormiga se moverá desde la ciudad i a la ciudad j según la probabilidad de transción y tras n iteraciones todas las hormigas habrán completado un ciclo y sus listas tabú contendrán un recorrido completo.
En este punto, para cada hormiga k se actualiza $L_k$ y la intensidad de feromona en cada arista del grafo según la fórmula dada anteriormente.
Además, se guarda el camino más corto encontrado, $min_k L_k, \quad k = 1,...,m$.
Este proceso se repite hasta que se cumple un criterio de parada, como alcanzar un número máximo de ciclos, $NC_{MAX}$ o todas las hormigas convergen a la misma solución.

Formalmente, el algoritmo se puede describir con los siguientes pasos:
\begin{enumerate}
    \item Inicialización: \\
    \hspace*{0.25cm} Definimos t := 0. \\
    \hspace*{0.25cm} Definimos NC := 0. \\
    \hspace*{0.25cm} Para cada arco (i, j) definimos la intensidad $\tau_{ij}(0) = c$ y $\Delta \tau_{ij} = 0$. \\
    \hspace*{0.25cm} Colocar las m hormigas en los n nodos.
    \item Sea s := 1. \\
    Para k = 1 hasta m hacer \\
    \hspace*{0.25cm} Colocar la ciudad inicial de la hormiga k en la lista tabú, $tabu_k(s)$. 

    \item Repetir hasta que la lista tabú se complete: \\
    Definimos s := s+1 \\
    Para k = 1 hasta m hacer \\
    \hspace*{0.25cm} Elegir la ciudad j a la que moverse con probabilidad $p_{ij}^k(t)$. \\
    \hspace*{0.25cm} Mover la hormiga k a la ciudad j. \\
    \hspace*{0.25cm} Colocar la ciudad j en la lista tabú, $tabu_k(s)$. 
    
    \item Para k = 1 hasta m hacer \\
    \hspace*{0.25cm} Mover la hormiga k de $tabu_k(n)$ a $tabu_k(1)$. \\
    \hspace*{0.25cm} Calcular la longitud del recorrido descrito por la hormiga k, $L_k$. \\
    \hspace*{0.25cm} Actualizar el camino más corto encontrado. \\
    \hspace*{0.25cm} Para cada arco (i, j) del grafo hacer \\
    \hspace*{0.25cm} Para k = 1 hasta m hacer 
    \begin{equation*}
        \Delta \tau_{ij}^k = \begin{cases}
        Q / L_k & \text{si } (i, j) \in \text{ciclo descrito por } tabu_k \\
        0 & \text{en otro caso}
        \end{cases}
    \end{equation*}

    \hspace*{0.25cm} $\Delta \tau_{ij} := \Delta \tau_{ij} + \Delta \tau_{ij}^{k}$

    \item Para cada arco (i, j) calcular $\tau_{ij}(t+n) = \rho \cdot \tau_{ij}(t) + \Delta \tau_{ij}$. \\
    Definir t := t + n. \\
    Definir NC := NC + 1. \\
    Para cada arco (i, j) definir $\Delta \tau_{ij} := 0$. 

    \item Si $(NC < NC_{MAX})$ y no todas las hormigas han convergido a la misma solución \\
    entonces \\
    \hspace*{0.25cm} Vaciar todas las listas tabú. \\
    \hspace*{0.25cm} Ir al paso 2. \\
    sino \\
    \hspace*{0.25cm} Imprimir el ciclo más corto encontrado. \\
    \hspace*{0.25cm} Parar.

\end{enumerate}

\subsection{Estrategia elitista}

Para mejorar el rendimiento del algoritmo, implementamos una estrategia elitista. 
Para ello añadimos al camino de cada arco del mejor ciclo la cantidad $e \cdot Q/L^{*}$, donde e es el número de hormigas elitistas y $L^{*}$ la longitud del mejor ciclo encontrado hasta el momento.
De esta forma, se refuerza el camino del mejor ciclo, acelerando la convergencia del algoritmo.

\section{Implementación}

\section{Conclusiones}

\section{Referencias}

\bibliographystyle{plain}
\bibliography{referencias} 

\end{document}